% ==============================================================================
% Chapter 3: Financial Ratios and Cash Flow Analysis
% Unit III of Pondicherry University MSQF 532 Curriculum
% ==============================================================================
\chapter{Financial Ratios and Cash Flow Analysis}

% ==============================================================================
% 3.1 Fundamentals of Ratio Analysis & Classification
% ==============================================================================
\section{Fundamentals of Ratio Analysis \& Classification}
\lecturedate{12/10/2026}

\begin{definition}{Financial Ratio \& Ratio Analysis}{def_financial_ratio}
A \textbf{Financial Ratio} is a quantitative mathematical expression establishing a proportional, fractional, or percentage relationship between two or more interrelated accounting figures extracted from corporate balance sheets, income statements, or cash flow reports.

\textbf{Ratio Analysis} is the diagnostic technique of standardizing and interpreting these numerical proportions to assess the operational efficiency, liquidity, leverage, solvency, and profitability of a business concern.
\end{definition}

\subsection{Managerial Purposes \& Core Advantages of Ratio Analysis}
Financial statement ratios serve seven strategic corporate functions:
\begin{enumerate}[leftmargin=*, topsep=3pt, itemsep=3pt]
    \item \textbf{Financial Forecasting \& Budgeting:} Historical ratio trajectories reveal structural growth rates, operating expense ratios, and asset utilization rhythms, facilitating predictive multi-year pro forma cash flow modeling.
    \item \textbf{Managerial Control:} Acts as a continuous diagnostic monitoring instrument by highlighting negative variances between actual operational performance and benchmark targets.
    \item \textbf{Facilitating Corporate Communication:} Distills complex, voluminous accounting transactions into standardized, intuitive indicator scores readily comprehended by board members, institutional investors, and lenders.
    \item \textbf{Measuring Operational Efficiency:} Quantifies the rotation velocity and conversion efficiency with which assets and working capital are transformed into revenues.
    \item \textbf{Investment Decision Support:} Equips equity analysts and shareholders with normalized valuation multiples (e.g. $\text{P/E}$, dividend yields, return on capital) to identify mispriced securities.
    \item \textbf{Evaluating Financial Solvency:} Simultaneously examines short-term debt-servicing capacity (liquidity) and long-term capital structure stability (financial gearing and leverage).
    \item \textbf{Inter-Firm \& Intra-Firm Benchmarking:} Facilitates cross-sectional peer comparison against industry leaders as well as longitudinal trend analysis across consecutive operating cycles.
\end{enumerate}

\subsection{Functional Taxonomy: Classification of Ratios by Purpose}
In corporate accounting, ratios are classified into three primary functional families based on the economic dimensions they diagnose:

\begin{figure}[htbp]
\centering
\begin{tikzpicture}[
    >=Stealth,
    banner/.style={
        draw=navyblue,
        fill=navyblue!12,
        text=navyblue,
        font=\bfseries\sffamily\small,
        rounded corners=2mm,
        text width=15.0cm,
        align=center,
        inner sep=7pt,
        line width=1.1pt
    },
    pillar/.style={
        draw=#1,
        fill=#1!4,
        rounded corners=2mm,
        text width=4.4cm,
        minimum height=5.4cm,
        align=flush left,
        inner sep=6pt,
        line width=1.1pt
    },
    line/.style={
        draw=navyblue,
        line width=1.2pt
    }
]

% Top Process Header
\node[banner] (top) {Classification of Financial Ratios by Purpose / Function};

% Middle Pillar (Turnover / Activity Ratios)
\node[pillar=amber, anchor=north] (p2) at ($(top.south) + (0, -1.3cm)$) {
    {\centering\color{amber}\textbf{\sffamily Turnover / Activity Ratios}\par}
    \vspace{0.3em}{\color{amber}\hrule height 0.6pt}\vspace{0.4em}
    \scriptsize
    \begin{itemize}[leftmargin=*, noitemsep, topsep=1pt]
        \item Stock / Inventory Turnover
        \item Debtors / Receivables Turnover
        \item Creditors / Payables Turnover
        \item Working Capital Turnover
        \item Fixed Assets Turnover
    \end{itemize}
    \vfill
    {\raggedright\tiny\color{amber}\textit{Measures operational velocity \& rotation efficiency of capital.}}
};

% Left Pillar (Profitability Ratios)
\node[pillar=darkteal, left=0.7cm of p2.north west, anchor=north east] (p1) {
    {\centering\color{darkteal}\textbf{\sffamily Profitability Ratios}\par}
    \vspace{0.3em}{\color{darkteal}\hrule height 0.6pt}\vspace{0.4em}
    \scriptsize
    \begin{itemize}[leftmargin=*, noitemsep, topsep=1pt]
        \item Return on Investment (ROI)
        \item Net Profit Ratio
        \item Gross Profit Ratio
        \item Operating Expenses Ratio
        \item Operating Profit Ratio
    \end{itemize}
    \vfill
    {\raggedright\tiny\color{darkteal}\textit{Diagnoses margin generation on revenue \& capital bases.}}
};

% Right Pillar (Solvency / Financial Ratios)
\node[pillar=navyblue, right=0.7cm of p2.north east, anchor=north west] (p3) {
    {\centering\color{navyblue}\textbf{\sffamily Solvency / Financial Ratios}\par}
    \vspace{0.3em}{\color{navyblue}\hrule height 0.6pt}\vspace{0.4em}
    \scriptsize
    \textbf{1. Short-Term Solvency (Liquidity):}
    \begin{itemize}[leftmargin=*, noitemsep, topsep=0.5pt]
        \item Current Ratio
        \item Liquidity / Quick Ratio
        \item Cash Position Ratio
    \end{itemize}
    \vspace{0.3em}
    \textbf{2. Long-Term Solvency (Leverage):}
    \begin{itemize}[leftmargin=*, noitemsep, topsep=0.5pt]
        \item Proprietary Ratio
        \item Debt-to-Equity Ratio
        \item Fixed Assets Ratio
        \item Capital Gearing Ratio
    \end{itemize}
};

% Coordinate anchor for horizontal branch bus (0.65cm below top banner)
\coordinate (bus) at ($(top.south) + (0, -0.65cm)$);

% Symmetrical Connector Bus & Downward Drop Stems (0.65cm stem clearance)
\draw[line] (top.south) -- (bus);
\draw[line] ($(p1.north |- bus)$) -- ($(p3.north |- bus)$);
\draw[line, ->] ($(p1.north |- bus)$) -- (p1.north);
\draw[line, ->] (bus) -- (p2.north);
\draw[line, ->] ($(p3.north |- bus)$) -- (p3.north);

\end{tikzpicture}
\caption{Taxonomy of Financial Ratios by Managerial Purpose and Analytical Function}
\label{fig:ratio_classification_tree}
\end{figure}

\subsection{Profitability Ratios: Formulation \& Operational Margin Identities}
\textbf{Profitability Ratios} measure an enterprise's ability to generate operational surplus, gross margin, and net earnings relative to revenue, total assets, or equity investment:

\begin{formulabox}[Profitability Ratios Formulations]
\begin{itemize}[leftmargin=*, noitemsep]
    \item \textbf{Gross Profit Ratio:} Reflects margin after direct production costs ($\COGS$):
    \begin{equation}
        \text{Gross Profit Ratio} = \frac{\text{Gross Profit}}{\text{Net Sales}} \times 100
    \end{equation}
    where $\text{Net Sales} = \text{Total Sales} - \text{Sales Returns}$, and $\text{Gross Profit} = \text{Net Sales} - \COGS$.

    \item \textbf{Operating Profit Ratio:} Measures operational profit after administrative, selling, and depreciation expenses:
    \begin{equation}
        \text{Operating Profit Ratio} = \frac{\text{Operating Profit}}{\text{Net Sales}} \times 100
    \end{equation}
    where $\text{Operating Profit} = \text{Gross Profit} - \text{Operating Expenses} = \text{Net Operating Income}$.

    \item \textbf{Operating Ratio:} Measures aggregate cost efficiency (direct + operating expenses per rupee of sales):
    \begin{equation}
        \text{Operating Ratio} = \frac{\COGS + \text{Operating Expenses}}{\text{Net Sales}} \times 100
    \end{equation}

    \item \textbf{Fundamental Invariance Identity:}
    \begin{equation}
        \text{Operating Ratio} + \text{Operating Profit Ratio} = 100\%
    \end{equation}

    \item \textbf{Net Profit Ratio:} Final margin available to equity holders after non-operating expenses and taxes:
    \begin{equation}
        \text{Net Profit Ratio} = \frac{\text{Net Profit After Tax } (\NPAT)}{\text{Net Sales}} \times 100
    \end{equation}
\end{itemize}
\end{formulabox}

\lecturedate{12/10/2026}
\begin{example}{Comprehensive Evaluation of Profitability Ratios}{ex_profitability_ratios}
\textbf{Problem:} \\
The following financial statement items relate to the trading and operating activities of a company for the year ended 31/12/1987:
\begin{center}
\small
\begin{tabular}{lr lr}
\toprule
\textbf{Particulars} & \textbf{Amount (\inr)} & \textbf{Particulars} & \textbf{Amount (\inr)} \\
\midrule
Sales & 1,00,000 & Salary of Salesmen & 1,800 \\
Purchases & 70,000 & Advertising & 700 \\
Closing Stock & 14,000 & Travelling Expenses & 500 \\
Sales Returns & 4,000 & Salaries (Office) & 3,000 \\
Dividend Received & 1,200 & Rent & 6,000 \\
Profit on Sale of Fixed Assets & 600 & Stationery & 200 \\
Loss on Sale of Shares & 300 & Depreciation & 1,000 \\
Opening Stock & 11,000 & Other Expenses & 2,000 \\
& & Provision for Tax & 7,000 \\
\bottomrule
\end{tabular}
\end{center}
You are required to prepare the formal Trading and Operating Profit \& Loss Account for the period, and compute:
\begin{enumerate*}[label=(\roman*), itemjoin=\qquad]
    \item Gross Profit Ratio,
    \item Operating Profit Ratio,
    \item Operating Ratio, and
    \item Net Profit Ratio.
\end{enumerate*}

\tcbline
\newpage
\textbf{Trading, Operating, and Profit \& Loss Account for the Year Ended 31/12/1987:}
\begin{center}
\small
\renewcommand{\arraystretch}{1.15}
\begin{tabularx}{\linewidth}{>{\raggedright\arraybackslash}X r >{\raggedright\arraybackslash}X r r}
\toprule
\textbf{Dr. Particulars} & \textbf{Amount (\inr)} & \textbf{Cr. Particulars} & \textbf{Inner (\inr)} & \textbf{Amount (\inr)} \\
\midrule
To Opening Stock & 11,000 & By Sales & 1,00,000 & \\
To Purchases & 70,000 & \textit{Less:} Sales Returns & (4,000) & 96,000 \\
To Gross Profit c/d & 29,000 & By Closing Stock & & 14,000 \\
\midrule
\textbf{Total} & \textbf{1,10,000} & \textbf{Total} & & \textbf{1,10,000} \\
\midrule
To Office Salaries & 3,000 & By Gross Profit b/d & & 29,000 \\
To Rent & 6,000 & & & \\
To Stationery & 200 & & & \\
To Depreciation & 1,000 & & & \\
To Other Operating Expenses & 2,000 & & & \\
To Salesmen Salary & 1,800 & & & \\
To Advertising & 700 & & & \\
To Travelling Expenses & 500 & & & \\
To Operating Profit c/d & \textbf{13,800} & & & \\
\midrule
\textbf{Total} & \textbf{29,000} & \textbf{Total} & & \textbf{29,000} \\
\midrule
To Loss on Sale of Shares & 300 & By Operating Profit b/d & & 13,800 \\
To Provision for Tax & 7,000 & By Dividend Received & & 1,200 \\
To Net Profit transferred to Equity & \textbf{8,300} & By Profit on Sale of Assets & & 600 \\
\midrule
\textbf{Total} & \textbf{15,600} & \textbf{Total} & & \textbf{15,600} \\
\bottomrule
\end{tabularx}
\end{center}

\textbf{Calculations of Profitability Ratios:}
\begin{enumerate}[leftmargin=*]
    \item \textbf{Gross Profit Ratio:}
    \begin{align*}
        \text{Gross Profit Ratio} &= \frac{\text{Gross Profit}}{\text{Net Sales}} \times 100 = \frac{29{,}000}{96{,}000} \times 100 = \mathbf{30.21\%}
    \end{align*}
    \textit{Interpretation:} The enterprise retains $\inr 30.21$ from every $\inr 100$ of net revenue after absorbing direct manufacturing and procurement costs ($\COGS$).

    \item \textbf{Operating Profit Ratio:}
    \begin{align*}
        \text{Operating Profit Ratio} &= \frac{\text{Operating Profit}}{\text{Net Sales}} \times 100 = \frac{13{,}800}{96{,}000} \times 100 = \mathbf{14.38\%}
    \end{align*}
    \textit{Interpretation:} Pure core operational business activities yield an operating margin of $14.38\%$ before considering capital gains, dividend windfalls, or income taxes.

    \item \textbf{Operating Ratio:}
    \begin{align*}
        \COGS &= \text{Opening Stock} + \text{Purchases} - \text{Closing Stock} \\
              &= 11{,}000 + 70{,}000 - 14{,}000 = \inr 67{,}000 \\
        \text{Operating Expenses} &= 3{,}000 + 6{,}000 + 200 + 1{,}000 + 2{,}000 + 1{,}800 + 700 + 500 \\
                                  &= \inr 15{,}200 \\
        \text{Total Operating Cost} &= \COGS + \text{Operating Expenses} = 67{,}000 + 15{,}200 = \inr 82{,}200 \\[0.3em]
        \text{Operating Ratio} &= \frac{\text{Total Operating Cost}}{\text{Net Sales}} \times 100 = \frac{82{,}200}{96{,}000} \times 100 = \mathbf{85.63\%}
    \end{align*}
    \textit{Verification of Invariance Identity:}
    \begin{equation*}
        \text{Operating Ratio} + \text{Operating Profit Ratio} = 85.63\% + 14.38\% = \mathbf{100.01\% \approx 100\%}
    \end{equation*}

    \item \textbf{Net Profit Ratio:}
    \begin{align*}
        \text{Net Profit Ratio} &= \frac{\text{Net Profit After Tax}}{\text{Net Sales}} \times 100 = \frac{8{,}300}{96{,}000} \times 100 = \mathbf{8.65\%}
    \end{align*}
    (If computed prior to non-operating adjustments as pure net operating income after tax: $\frac{13,800 - 7,000}{96,000} \times 100 = \frac{6,800}{96,000} \times 100 = 7.08\%$).
\end{enumerate}

\textbf{Final Results Summary:}
\begin{itemize}[noitemsep]
    \item Gross Profit Ratio = $\mathbf{30.21\%}$
    \item Operating Profit Ratio = $\mathbf{14.38\%}$
    \item Operating Ratio = $\mathbf{85.63\%}$
    \item Net Profit Ratio = $\mathbf{8.65\%}$ (or $\mathbf{7.81\%}$ based on standard pre-dividend adjusted net line)
\end{itemize}
\end{example}

% ==============================================================================
% 3.2 Turnover / Activity Ratios & Operational Velocity
% ==============================================================================
\section{Turnover / Activity Ratios \& Operational Velocity}
\lecturedate{13/10/2026}

\begin{definition}{Turnover Ratio \& Operational Velocity}{def_turnover_ratio}
\textbf{Turnover Ratios} (also termed \textbf{Activity Ratios} or \textbf{Performance Ratios}) measure the operational efficiency and rotation velocity with which an enterprise converts its balance sheet assets and capital resources into revenue or cost of sales.

Unlike profitability metrics expressed as percentages, turnover ratios are mathematically expressed in \textbf{number of times} (rotation frequency per annum). A higher multiple generally signals superior managerial agility, efficient working capital deployment, and minimized idle capital.
\end{definition}

\begin{formulabox}[Comprehensive Activity and Velocity Formulations]
\begin{itemize}[leftmargin=*, noitemsep]
    \item \textbf{Inventory / Stock Turnover Ratio (ITR):}
    \begin{equation}
        \text{ITR} = \frac{\text{Cost of Goods Sold } (\COGS)}{\text{Average Inventory}} \quad \text{or} \quad \frac{\text{Net Sales}}{\text{Average Inventory}}
    \end{equation}
    where $\text{Average Inventory} = \dfrac{\text{Opening Stock} + \text{Closing Stock}}{2}$.
    \begin{equation}
        \text{Inventory Holding Period (Velocity)} = \frac{365}{\text{ITR}} \text{ days} = \frac{12}{\text{ITR}} \text{ months}
    \end{equation}

    \item \textbf{Debtors / Accounts Receivable Turnover Ratio (DTR):}
    \begin{equation}
        \text{DTR} = \frac{\text{Net Credit Sales}}{\text{Average Accounts Receivable}}
    \end{equation}
    where $\text{Net Credit Sales} = \text{Credit Sales} - \text{Sales Returns}$, and $\text{Average Receivables} = \text{Debtors} + \text{Bills Receivable}$.
    \begin{equation}
        \text{Debtors Collection Period (Velocity)} = \frac{365}{\text{DTR}} \text{ days} = \frac{12}{\text{DTR}} \text{ months}
    \end{equation}

    \item \textbf{Creditors / Accounts Payable Turnover Ratio (CTR):}
    \begin{equation}
        \text{CTR} = \frac{\text{Net Credit Purchases}}{\text{Average Accounts Payable}}
    \end{equation}
    where $\text{Net Credit Purchases} = \text{Credit Purchases} - \text{Purchase Returns}$, and $\text{Average Payables} = \text{Creditors} + \text{Bills Payable}$.
    \begin{equation}
        \text{Creditors Payment Period (Velocity)} = \frac{365}{\text{CTR}} \text{ days} = \frac{12}{\text{CTR}} \text{ months}
    \end{equation}

    \item \textbf{Working Capital Turnover Ratio:}
    \begin{equation}
        \text{Working Capital Turnover} = \frac{\text{Net Sales (or Cost of Sales)}}{\text{Net Working Capital}}
    \end{equation}
    where $\text{Net Working Capital} = \text{Current Assets} - \text{Current Liabilities}$.

    \item \textbf{Fixed Assets Turnover Ratio:}
    \begin{equation}
        \text{Fixed Assets Turnover} = \frac{\COGS}{\text{Net Fixed Assets}} \quad \text{or} \quad \frac{\text{Net Sales}}{\text{Net Fixed Assets}}
    \end{equation}

    \item \textbf{Capital Turnover Ratio:}
    \begin{equation}
        \text{Capital Turnover} = \frac{\text{Cost of Sales (or Net Sales)}}{\text{Capital Employed}}
    \end{equation}
\end{itemize}
\end{formulabox}

\lecturedate{13/10/2026}
\begin{example}{Comprehensive Calculation of Turnover Ratios and Velocity Periods}{ex_turnover_comprehensive}
\textbf{Problem:} \\
The accounting records of a manufacturing corporation disclose the following operating data for the financial year:
\begin{center}
\small
\begin{tabularx}{\linewidth}{>{\raggedright\arraybackslash}X r >{\raggedright\arraybackslash}X r}
\toprule
\textbf{Particulars} & \textbf{Amount (\inr)} & \textbf{Particulars} & \textbf{Amount (\inr)} \\
\midrule
Total Sales & 5,00,000 & Cash Purchases & 50,000 \\
Cash Sales & 40,000 & Purchase Returns & 10,000 \\
Sales Returns & 20,000 & Creditors at the end & 60,000 \\
Debtors at the end & 80,000 & Bills Payable at the end & 20,000 \\
Bills Receivable at the end & 20,000 & Reserve on Discount on Creditors & 2,000 \\
Reserve for Doubtful Debts & 1,000 & Opening Stock & 50,000 \\
Total Purchases & 3,00,000 & Closing Stock & 40,000 \\
Gross Profit & 1,00,000 & Fixed Assets & 10,00,000 \\
\bottomrule
\end{tabularx}
\end{center}
From the above information, compute:
\begin{enumerate*}[label=(\alph*), itemjoin=\quad]
    \item Stock Turnover Ratio,
    \item Fixed Asset Turnover Ratio (under both Cost of Sales and Net Sales bases),
    \item Debtors Turnover Ratio and Debtors Collection Period (in months and days),
    \item Creditors Turnover Ratio and Creditors Payment Period (in months and days).
\end{enumerate*}

\tcbline
\textbf{Reconciliation of Intermediate Financial Aggregates:}
\begin{enumerate}[leftmargin=*]
    \item \textbf{Net Sales:}
    \begin{equation*}
        \text{Net Sales} = \text{Total Sales} - \text{Sales Returns} = 5{,}00{,}000 - 20{,}000 = \mathbf{\inr 4{,}80{,}000}
    \end{equation*}

    \item \textbf{Cost of Goods Sold (Cost of Sales):}
    \begin{equation*}
        \COGS = \text{Net Sales} - \text{Gross Profit} = 4{,}80{,}000 - 1{,}00{,}000 = \mathbf{\inr 3{,}80{,}000}
    \end{equation*}

    \item \textbf{Average Stock:}
    \begin{equation*}
        \text{Average Stock} = \frac{\text{Opening Stock} + \text{Closing Stock}}{2} = \frac{50{,}000 + 40{,}000}{2} = \mathbf{\inr 45{,}000}
    \end{equation*}

    \item \textbf{Net Credit Sales:}
    \begin{align*}
        \text{Net Credit Sales} &= \text{Total Sales} - \text{Cash Sales} - \text{Sales Returns} \\
                                &= 5{,}00{,}000 - 40{,}000 - 20{,}000 = \mathbf{\inr 4{,}40{,}000}
    \end{align*}

    \item \textbf{Average Accounts Receivable:}
    (Note: Reserves for doubtful debts are provisions and are not deducted from gross receivables in turnover velocity analysis):
    \begin{equation*}
        \text{Accounts Receivable} = \text{Debtors} + \text{Bills Receivable} = 80{,}000 + 20{,}000 = \mathbf{\inr 1{,}00{,}000}
    \end{equation*}

    \item \textbf{Net Credit Purchases:}
    \begin{align*}
        \text{Net Credit Purchases} &= \text{Total Purchases} - \text{Cash Purchases} - \text{Purchase Returns} \\
                                    &= 3{,}00{,}000 - 50{,}000 - 10{,}000 = \mathbf{\inr 2{,}40{,}000}
    \end{align*}

    \item \textbf{Average Accounts Payable:}
    \begin{equation*}
        \text{Accounts Payable} = \text{Creditors} + \text{Bills Payable} = 60{,}000 + 20{,}000 = \mathbf{\inr 80{,}000}
    \end{equation*}
\end{enumerate}

\tcbline
\textbf{Step-by-Step Ratio Computations:}
\begin{enumerate}[leftmargin=*]
    \item \textbf{Stock Turnover Ratio:}
    \begin{equation*}
        \text{Stock Turnover Ratio} = \frac{\COGS}{\text{Average Stock}} = \frac{3{,}80{,}000}{45{,}000} = \mathbf{8.44 \text{ times}}
    \end{equation*}
    (If evaluated on Net Sales basis: $\frac{4,80,000}{45,000} = 10.67 \text{ times}$).

    \item \textbf{Fixed Asset Turnover Ratio:}
    \begin{itemize}[noitemsep]
        \item \textit{Based on Cost of Sales:}
        \begin{equation*}
            \text{Fixed Asset Turnover} = \frac{\COGS}{\text{Fixed Assets}} = \frac{3{,}80{,}000}{10{,}00{,}000} = \mathbf{0.38 \text{ times}}
        \end{equation*}
        \item \textit{Based on Net Sales:}
        \begin{equation*}
            \text{Fixed Asset Turnover} = \frac{\text{Net Sales}}{\text{Fixed Assets}} = \frac{4{,}80{,}000}{10{,}00{,}000} = \mathbf{0.48 \text{ times}}
        \end{equation*}
    \end{itemize}

    \item \textbf{Debtors Turnover Ratio and Debtors Collection Velocity:}
    \begin{equation*}
        \text{Debtors Turnover Ratio} = \frac{\text{Net Credit Sales}}{\text{Accounts Receivable}} = \frac{4{,}40{,}000}{1{,}00{,}000} = \mathbf{4.40 \text{ times}}
    \end{equation*}
    \begin{align*}
        \text{Collection Period (Months)} &= \frac{12}{\text{Debtors Turnover}} = \frac{12}{4.40} = \mathbf{2.73 \text{ months}} \\[0.4em]
        \text{Collection Period (Days)} &= \frac{365}{\text{Debtors Turnover}} = \frac{365}{4.40} = \mathbf{82.95 \approx 83 \text{ days}}
    \end{align*}

    \item \textbf{Creditors Turnover Ratio and Creditors Payment Velocity:}
    \begin{equation*}
        \text{Creditors Turnover Ratio} = \frac{\text{Net Credit Purchases}}{\text{Accounts Payable}} = \frac{2{,}40{,}000}{80{,}000} = \mathbf{3.00 \text{ times}}
    \end{equation*}
    \begin{align*}
        \text{Payment Period (Months)} &= \frac{12}{\text{Creditors Turnover}} = \frac{12}{3.00} = \mathbf{4.00 \text{ months}} \\[0.4em]
        \text{Payment Period (Days)} &= \frac{365}{\text{Creditors Turnover}} = \frac{365}{3.00} = \mathbf{121.67 \approx 122 \text{ days}}
    \end{align*}
\end{enumerate}

\textbf{Final Results Summary:}
\begin{itemize}[noitemsep]
    \item Stock Turnover Ratio = $\mathbf{8.44 \text{ times}}$
    \item Fixed Asset Turnover Ratio = $\mathbf{0.38 \text{ times}}$ (Cost of Sales) / $\mathbf{0.48 \text{ times}}$ (Net Sales)
    \item Debtors Turnover Ratio = $\mathbf{4.40 \text{ times}}$ ($\mathbf{2.73 \text{ months}}$ or $\mathbf{83 \text{ days}}$)
    \item Creditors Turnover Ratio = $\mathbf{3.00 \text{ times}}$ ($\mathbf{4.00 \text{ months}}$ or $\mathbf{122 \text{ days}}$)
\end{itemize}
\end{example}
